\section{Canonical forms and coherent direct images}\label{sec:canonical-coherent}

The duality calculation above retains the holomorphic canonical line.
We now determine that line from the quotient charts. The same charts
also determine the relative logarithmic forms and their coherent
direct images. These calculations will give both the restriction map
on the Picard group and the ordinary Hodge groups.

Throughout this section, fix the sufficiently small nonzero parameter
$u$, integral lifts, logarithmic increments and adapted cusp fan of
Theorem~\ref{cmp:thm:comparison}. Write
\[
 f:Y=Y^{\mathrm{ad}}(u)\longrightarrow B=\mathbb P^1.
\]
Let $F_3,F_4$ be the reduced smooth fibers over $0,1$, and let
$W=Z_\infty$ be the full reduced toric fiber. In particular,
\begin{equation}\label{nf:eq:fiber-divisors}
 f^*(0)=3F_3,\qquad f^*(1)=4F_4,\qquad f^*(\infty)=W.
\end{equation}
The divisor $W$ has normal crossings locally, including the distinct
branches at its self-intersections. All coherent sheaves below are
analytic sheaves over $\mathbb C$.

\subsection{The canonical divisor}

The invariant elliptic differential and the logarithmic differential
along the divisor line give a top form before taking the quotient.
Their wedge removes the ambiguity in choosing a line frame.

\begin{theorem}\label{nf:thm:canonical}
The canonical line of the selected threefold is
\begin{equation}\label{nf:eq:canonical}
 K_Y\simeq\mathcal O_Y(2F_4-W)
     \simeq f^*\mathcal O_B(-1)\otimes\mathcal O_Y(2F_4)
     \simeq\mathcal O_Y(-2F_4).
\end{equation}
For every integer $r>0$, one has $H^0(Y,K_Y^{\otimes r})=0$.
\end{theorem}

\begin{proof}
Over $B^\circ=B\setminus\{0,1,\infty\}$, let
$\alpha_t=dx/(2y)$ be the invariant relative differential on the
elliptic surface. In a local frame of $M=\mathcal O(P-O)$, write a
nonzero line vector as $v$. Set
\begin{equation}\label{nf:eq:regular-volume}
 \Omega=dt\wedge\alpha_t\wedge\frac{dv}{v}.
\end{equation}
An absolute lift of $\alpha_t$ is understood in this expression.
Changing it adds a multiple of $dt$ and leaves $\Omega$ unchanged.
If $v_j=g_{ij}(x,t)v_i$, the change in $dv/v$ is $d\log g_{ij}$.
This form comes from the elliptic surface, so it wedges to zero
with its top form $dt\wedge\alpha_t$. Thus $\Omega$ is a global
nowhere-zero form on $M^\times$ over $B^\circ$.

The actual contraction generator is a fiber-linear lift of
translation by $-6P$. Locally it has the form
\[
 (x,v)\longmapsto(t_{-6P}x,G(x,t)v).
\]
Translation preserves $\alpha_t$ modulo base differentials.
The added $d\log G$ again comes from the elliptic surface.
Consequently the generator, including its fixed scalar $u$,
preserves \eqref{nf:eq:regular-volume}. The form descends to the
regular family of two-tori.

On each smooth root family, put
\[
 \alpha_i=\frac{dX}{2Y},\qquad
 \sigma_i=ds\wedge\alpha_i\wedge\frac{dv}{v}
 \quad(i=0,1).
\]
The same frame and translation argument makes $\sigma_i$ a
nowhere-zero holomorphic form on the proper smooth torus family.
The transported meromorphic line identification over the punctured
root disk can multiply $v$ by a function of the surface coordinates
and a power of $s$. Its logarithmic derivative does not change this
top wedge.

The two actual substitutions give
\begin{align}
 t=s^3,\quad x=s^4X,\quad y=s^6Y
 &\quad\Longrightarrow\quad
 \Omega=3\sigma_0, \label{nf:eq:volume-three}\\
 t=1+s^4,\quad x=s^2X,\quad y=s^3Y
 &\quad\Longrightarrow\quad
 \Omega=4s^2\sigma_1. \label{nf:eq:volume-four}
\end{align}
The reduced equations are respectively
\[
 Y^2=X^3-3s(s^3-1)X+2(s^3-1)^2
\]
and
\[
 Y^2=X^3-3(1+s^4)^3X+2s^2(1+s^4)^4.
\]
Their central elliptic fibers are smooth. The vanishing orders in
\eqref{nf:eq:volume-three}--\eqref{nf:eq:volume-four} are therefore
exactly $0$ and $2$.

The punctured gluing maps are translations by the specified
logarithmic torus sections. A translation preserves the invariant
relative two-form; its base derivative disappears after wedging
with $ds$. Hence these gluing maps preserve both displayed absolute
forms. The finite affine actions are free on the total spaces, so
their quotient maps are \'{e}tale there. The orders $0$ and $2$
descend unchanged to $F_3$ and $F_4$. The ramification of the base
maps does not divide these divisor orders.

The forward coordinate decks also fix the signs directly. At order
three, $(s,X,Y)\mapsto(\zeta s,\zeta^2X,Y)$ gives
$\sigma_0\mapsto\zeta^3\sigma_0=\sigma_0$. At order four,
$(s,X,Y)\mapsto(is,-X,iY)$ gives
\begin{equation}\label{nf:eq:volume-character}
 g^*\sigma_1=-\sigma_1,\qquad
 g^*(s^2\sigma_1)=s^2\sigma_1.
\end{equation}
The fiber-linear multiplier adds only terms already killed by the
top wedge, and the affine translations preserve relative volume.
In particular omitting $s^2$ would leave
$g^*\sigma_1-\sigma_1=-2\sigma_1$.

At infinity use $q=1/t$, $x_\infty=q^2x$ and $y_\infty=q^3y$.
Modulo $dq$,
\[
 \alpha_t=q\alpha_\infty,\qquad
 \alpha_\infty=\frac{dx_\infty}{2y_\infty},\qquad
 dt=-q^{-2}dq.
\]
In the actual Tate coordinate $z$ and divisor-line coordinate $w$,
this gives
\begin{equation}\label{nf:eq:cusp-volume}
 \Omega=-q^{-1}dq\wedge\alpha_\infty\wedge d\log w,\qquad
 \alpha_\infty=a(q)d\log z,
\end{equation}
where $a$ is a holomorphic unit. To see the unit assertion, use the
whole-disk Tate comparison in the construction. Tate's functions
satisfy $z\,dx(z)/dz=x(z)+2y(z)$, so their invariant differential is
$dx/(2y+x)=d\log z$ \cite{TateElliptic}. Completing the square and
shifting $x$ preserve that differential. The short Weierstrass
coefficients at the central fiber are $(-1/48,1/864)$; those of the
present equation are $(-3,2)$. Their holomorphic scaling is a unit,
with square $12$ at zero. Thus
$a(0)=\pm1/(2\sqrt3)$, in particular $a(0)\ne0$.

The central cubic gives the same calibration without a choice of
period logarithms. It is
$y_\infty^2=(x_\infty-1)^2(x_\infty+2)$. With
$\eta=y_\infty/(x_\infty-1)$,
\begin{equation}\label{nf:eq:nodal-differential}
 x_\infty=\eta^2-2,\quad
 y_\infty=\eta(\eta^2-3),\quad
 \alpha_\infty=\frac{d\eta}{\eta^2-3}
 =\frac1{2\sqrt3}\,
 d\log\frac{\eta-\sqrt3}{\eta+\sqrt3}.
\end{equation}
Its nodal residues are nonzero.

On a maximal chart of the height-one unimodular cusp fan,
$q=x_0x_1x_2$. The character-coordinate change has determinant
$\pm1$, and therefore
\begin{equation}\label{nf:eq:toric-volume}
 dq\wedge d\log z\wedge d\log w
       =\pm dx_0\wedge dx_1\wedge dx_2.
\end{equation}
Period multipliers and the unit translation of the retained origin
add only base differentials. The adapted shear
$(z,w)\mapsto(z/w^2,w)$ has determinant one on logarithmic
differentials. Thus \eqref{nf:eq:toric-volume} applies to the chosen
adapted quotient, with all its units retained. Its free period
action is locally biholomorphic. Equation~\eqref{nf:eq:cusp-volume}
has a simple pole on every local branch of $W$.

All local expressions came from the same form on the regular
family and agree on the actual annular overlaps. They give a
meromorphic section of $K_Y$ with divisor $2F_4-W$.
Since $W\sim4F_4$ by \eqref{nf:eq:fiber-divisors}, this proves
\eqref{nf:eq:canonical}. Finally,
$\mathcal O_Y(-2rF_4)$ is an ideal subsheaf of $\mathcal O_Y$.
Every holomorphic function on compact connected $Y$ is constant,
and a constant vanishing along the nonempty divisor $F_4$ is zero.
This proves the pluricanonical vanishing.
\end{proof}

\subsection{Relative logarithmic forms and the scalar action}

Put
\[
 \Sigma_B=0+1+\infty,\qquad D=F_3+F_4+W,\qquad
 \mathscr L_B=\Omega_B^1(\log\Sigma_B)\simeq\mathcal O_B(1),
\]
and define
\begin{equation}\label{nf:eq:relative-forms}
 \mathscr E=
 \Omega_Y^1(\log D)/f^*\mathscr L_B,\qquad
 \mathcal H=f_*\mathscr E.
\end{equation}
At a self-intersection of $W$, logarithmic forms use the separate
local normal-crossing branches. No equality of their residues is
part of this definition. On a finite root cover, the base form is
$m\,ds/s$ and its quotient removes the logarithmic normal
direction. At a cusp chart it is
$d\log q=d\log x_0+d\log x_1+d\log x_2$.
These local descriptions give an exact sequence of vector bundles
\begin{equation}\label{nf:eq:absolute-log-sequence}
 0\longrightarrow f^*\mathscr L_B
 \longrightarrow\Omega_Y^1(\log D)
 \longrightarrow\mathscr E\longrightarrow0,
\end{equation}
with $\mathscr E$ of rank two.

Scalar multiplication on the original line $M$ commutes with its
fiber-linear lift. Let $V$ be the resulting infinitesimal vector,
normalized by $V(v)=v$. It descends through the regular quotient,
the finite affine actions and the logarithmic translations.
In the original cusp coordinates $V=w\partial_w$. In the adapted
coordinates $\zeta_1=z/w^2$, $\zeta_2=w$, it is
\begin{equation}\label{nf:eq:scalar-vector}
 V=-2\zeta_1\partial_{\zeta_1}+\zeta_2\partial_{\zeta_2}.
\end{equation}
This is nonzero in the relative logarithmic tangent bundle,
including at toric fixed points. Its ordinary tangent value at
such a point need not be nonzero. Contraction gives a surjection
of bundles on $Y$; denote its kernel by $\mathcal K_{\mathrm{rel}}$:
\begin{equation}\label{nf:eq:scalar-sequence}
 0\longrightarrow\mathcal K_{\mathrm{rel}}
 \longrightarrow\mathscr E\xrightarrow{\ \iota_V\ }
 \mathcal O_Y\longrightarrow0.
\end{equation}

\begin{theorem}\label{nf:thm:relative-line}
The direct image $\mathcal H$ is locally free of rank two.
Contraction induces an exact sequence
\begin{equation}\label{nf:eq:H-sequence}
 0\longrightarrow\mathcal L
 \longrightarrow\mathcal H\xrightarrow{\ q_V\ }\mathcal O_B
 \longrightarrow0,\qquad
 \mathcal L=\mathcal O_B(\infty-0-1)\simeq\mathcal O_B(-1).
\end{equation}
Its splitting as a right inverse to $q_V$ is unique.
Equivalently, there is a unique global relative logarithmic form
$\theta$ satisfying $\theta(V)=1$.
\end{theorem}

\begin{proof}
The proper map has connected fibers and normal total space and
base. Stein factorization gives $f_*\mathcal O_Y=\mathcal O_B$:
its finite factor is birational to the normal curve $B$.
Pushforward of contraction therefore has the stated target.

On a small smooth torus-family disk, invariant relative
differentials form the dual of the relative Lie algebra.
The scalar vector is nonzero and has a holomorphic dual form
with value one. On either good-reduction root disk, average such
a form over the actual forward cyclic action, including its
action on the base. The average remains holomorphic and has value
one on $V$. Its kernel consists of the horizontal elliptic forms.
Exact finite pushforward and averaging descend this description
to the finite quotient.

At the cusp the relative classes of $d\log z,d\log w$ are a
frame on every height-one unimodular chart. Period changes add
multiples of $d\log q$, so these classes descend. Their
contractions are $0,1$. In the adapted coordinates they are
$d\log\zeta_1+2d\log\zeta_2$ and $d\log\zeta_2$.
Using $f_*\mathcal O_Y=\mathcal O_B$ now gives a frame of
$\mathcal H$ over the cusp disk. Thus $\mathcal H$ is locally
free, $q_V$ is onto, and its kernel is a line
$\mathcal L=f_*\mathcal K_{\mathrm{rel}}$.

The invariant elliptic differential
$\omega=dx/(2y)$ defines a meromorphic section of $\mathcal L$.
Elliptic translations and the selected logarithmic gluing preserve
it as a relative differential. Its divisor in this base line can
be calculated on the root disks. The forward deck characters and
invariant holomorphic generators are
\begin{equation}\label{nf:eq:horizontal-frames}
 \begin{array}{c|c|c|c}
 \text{base point}&\text{root coordinate}&g^*\alpha_i&
       \text{generator of }\mathcal L\\ \hline
 0&t=s^3&\zeta^2\alpha_0&s\alpha_0\\
 1&t-1=s^4&i\alpha_1&s^3\alpha_1 .
 \end{array}
\end{equation}
Consequently
\[
 \omega=t^{-1}(s\alpha_0),\qquad
 \omega=(t-1)^{-1}(s^3\alpha_1).
\]
Each pole has order one on the base. These are generators of a
direct-image line on the base; their pullbacks need not be
nonvanishing along the central root fiber.
The forward deck is essential here:
$g^*(s\alpha_1)-s\alpha_1=-2s\alpha_1$, whereas
$g^*(s^3\alpha_1)=s^3\alpha_1$.

At infinity, \eqref{nf:eq:cusp-volume} gives
$\omega=q\,a(q)d\log z$, with $a(0)\ne0$.
The horizontal logarithmic class $d\log z$ is a nonvanishing
frame there. Thus
\[
 \operatorname{div}_{\mathcal L}(\omega)=\infty-0-1.
\]
This identifies the actual line and its meromorphic frame in
\eqref{nf:eq:H-sequence}. The extension splits because
$H^1(B,\mathcal O_B(-1))=0$. Two right inverses differ by a
section of $\mathcal O_B(-1)$, whose $H^0$ is zero.
This proves uniqueness.
\end{proof}

The form $\theta$ is a section of the relative logarithmic
quotient. The theorem does not assert that it is an absolute
holomorphic one-form. It does split \eqref{nf:eq:scalar-sequence}
on $Y$, by sending $1$ to $\theta$.

\subsection{The coherent algebra of the cusp}

The central toric fiber is nonnormal. Its coherent cohomology is
therefore computed by its normalization and the identifications of
the boundary curves, rather than by the normalization alone.

Let $\nu:S\to W$ be the normalization. Here $S$ is the blowup of
$\mathbb P^2$ at its three torus fixed points. Label its boundary
hexagon $B_1,\ldots,B_6$ cyclically, with $B_i$ paired with
$B_{i+3}$. Let $C_1,C_2,C_3$ be their reduced image double curves
and $p_+,p_-$ the two triple points.

\begin{lemma}\label{nf:lem:cusp-coherent}
One has
\[
 h^q(W,\mathcal O_W)=(1,2,1)\quad(q=0,1,2).
\]
The natural comparison $H^1(W;\mathbb C)\to H^1(W,\mathcal O_W)$
is an isomorphism. In degree two it restricts to an isomorphism
\[
 \ker\!\left(H^2(W;\mathbb C)\xrightarrow{\nu^*}
                  H^2(S;\mathbb C)\right)
       \xrightarrow{\ \sim\ }H^2(W,\mathcal O_W).
\]
The cup product
$\bigwedge^2H^1(W,\mathcal O_W)\to H^2(W,\mathcal O_W)$
is an isomorphism.
\end{lemma}

\begin{proof}
Orient each pair by $+B_i,-B_{i+3}$. The branch orders at $p_+,p_-$
give the alternating restriction $g_1-g_2+g_3$ at both points.
The normalization sequence is
\begin{equation}\label{nf:eq:normalization-complex}
 0\longrightarrow\mathcal O_W
 \longrightarrow\nu_*\mathcal O_S
 \longrightarrow\bigoplus_{i=1}^3\mathcal O_{C_i}
 \longrightarrow\mathbb C_{p_+}\oplus\mathbb C_{p_-}
 \longrightarrow0.
\end{equation}
The middle arrow takes the differences of restrictions through
the actual holomorphic gluing maps. The unit factors can change
those maps without changing their branch incidence.
This is the normalization sequence also used in
\cite[Lemma~9.12]{TorusConstruction}.

Exactness is local. At a double point it is the fiber-product
sequence for $\mathbb C\{x,y,z\}/(xy)$. At a triple point it is
the alternating branch complex for
$\mathbb C\{x,y,z\}/(xyz)$. To verify the latter, suppose branch
functions $f_1(y,z),f_2(x,z),f_3(x,y)$ have common restrictions
$g_{12}(z),g_{13}(y),g_{23}(x)$ and common value $c$ at the origin.
They glue to the restriction of
\[
 f_1+f_2+f_3-g_{12}-g_{13}-g_{23}+c.
\]
For exactness at the double-curve term, suppose instead that
$g_{12}(0)-g_{13}(0)+g_{23}(0)=0$. Set
\[
 f_3=0,\quad
 f_1(y,z)=g_{13}(y)+g_{12}(z)-g_{12}(0),\quad
 f_2(x,z)=g_{23}(x).
\]
Their alternating differences are the prescribed $g_{ij}$.
The last map is onto. The same local simplicial complexes are
exact for constant sheaves, and inclusion of constants gives a
termwise comparison with \eqref{nf:eq:normalization-complex}.

The rational surface $S$ and the curves $C_i\simeq\mathbb P^1$
have no positive coherent cohomology. Hence the coherent complex
of global sections is
\begin{equation}\label{nf:eq:cusp-constants}
 \mathbb C\xrightarrow{\ 0\ }\mathbb C^3
 \xrightarrow{\ (a,b,c)\mapsto(a-b+c)(1,1)\ }\mathbb C^2.
\end{equation}
Its cohomology has dimensions $(1,2,1)$.
For constant coefficients, the degree-zero row of the
normalization spectral sequence is the same complex and the
degree-one row is zero. Thus the degree-one comparison is an
isomorphism. In total degree two, the kernel of restriction to
$S$ is its one-dimensional $(2,0)$ term. No differential enters
it from degree-one cohomology of the strata, and there is no
normalization term beyond degree two. This term maps
isomorphically to $H^2(W,\mathcal O_W)$.

It remains to identify a nonzero product in this term. With unit
periods, the angular-quotient description
\eqref{c26:eq:Z} projects to the polygon torus $H/\Gamma$.
Angular coordinate zero gives a section of that projection.
The product of the two polygon-period covectors is therefore
nonzero in $H^2(W;\mathbb C)$. Its restriction to $S$ is zero:
the map from the simply connected surface $S$ to the polygon
torus lifts to its contractible universal cover.
The smooth transport \eqref{c26:eq:transport} preserves the
strata and lifts to the normalization, so it carries this
topological product statement to the actual cusp units.
The coherent comparison is still the one defined by the actual
holomorphic sequence \eqref{nf:eq:normalization-complex}.
Inclusion of constants respects cup products, so the nonzero
product survives in $H^2(W,\mathcal O_W)$. Both its source and
target are one-dimensional.
\end{proof}

\subsection{The structure-sheaf direct images}

Write
\[
 \mathcal C_q=R^qf_*\mathcal O_Y,\qquad \mathcal G=\mathcal C_1.
\]
The primitive invariant cohomology covector in (1.9) is
$\psi=e_4^\vee$. We will denote its coherent image by $\eta$;
the two symbols distinguish the integral class from its
$(0,1)$ component.

\begin{theorem}\label{nf:thm:coherent}
The sheaves $\mathcal C_q$ vanish for $q>2$, and
\begin{equation}\label{nf:eq:coherent-bundles}
 \mathcal C_0=\mathcal O_B,\qquad
 \mathcal C_1\simeq\mathcal O_B\oplus\mathcal O_B(-1),\qquad
 \mathcal C_2\simeq\mathcal O_B(-1).
\end{equation}
The section $\eta=\psi^{0,1}$ defines a subbundle and the cup
product defines an isomorphism
\begin{equation}\label{nf:eq:G-sequence}
 0\longrightarrow\mathcal O_B\eta\longrightarrow\mathcal G
 \longrightarrow\mathcal Q\longrightarrow0,\qquad
 \mathcal Q\simeq\mathcal O_B(-1),\qquad
 \bigwedge^2\mathcal G\xrightarrow{\ \sim\ }\mathcal C_2.
\end{equation}
The first sequence splits, without a distinguished choice of
complement. In particular,
\[
 h^q(Y,\mathcal O_Y)=(1,1,0,0),\qquad
 h^q(Y,K_Y)=(0,0,1,1).
\]
\end{theorem}

\begin{proof}
First we establish local freeness on the full base. The map $f$
is flat: a nonzero germ on the smooth base acts injectively on
the regular local domain of $Y$, and torsion-free modules over
the base discrete valuation ring are flat.
Near infinity the scheme fiber is the reduced divisor $W$.
Its coherent dimensions from Lemma~\ref{nf:lem:cusp-coherent}
are those of the nearby two-tori. Grauert's proper analytic
cohomology and base-change theorem
\cite[\S7]{GrauertDirect}, with its correction
\cite{GrauertCorrection}, therefore gives
locally free direct images of ranks $(1,2,1)$, with their
fiber-cohomology comparison, across infinity.

At either finite point use the proper smooth root family
$g:J'\to\Delta_s$, the free affine cyclic action of order $m$,
and the base map $\pi(s)=s^m$. Exact finite pushforward and
averaging give
\begin{equation}\label{nf:eq:finite-coherent-descent}
 R^q(f|_N)_*\mathcal O_N
       \simeq\left(\pi_*R^qg_*\mathcal O_{J'}\right)^{\mu_m}.
\end{equation}
The bundles on the root disk are free of ranks $(1,2,1)$.
A finite semilinear action admits holomorphic eigenframes:
diagonalize its central fiber, lift that basis, and apply the
character projectors obtained by averaging. Their determinant
remains a unit after shrinking the disk. If a forward eigenframe
$e$ has weight $b$, meaning $g^*e=\zeta^b e$, its invariant
module is generated over $\mathcal O_{\Delta_t}$ by $s^re$,
where $0\le r<m$ and $r+b\equiv0\pmod m$.
Equation~\eqref{nf:eq:finite-coherent-descent} thus proves local
freeness at the multiple fibers. It computes the full quotient;
it does not replace the scheme fiber $mF_m$ by its reduction.
These local descriptions also prove the vanishing for $q>2$.

On $B^\circ$ the real integral flat class $\psi$ has a
holomorphic image
$\eta=\psi^{0,1}$ under
$R^1f_*\mathbb C\otimes\mathcal O_B\to\mathcal G$.
This image is nonzero on every smooth torus. Indeed, if a
nonzero real class had zero $(0,1)$ component, it would belong
both to $H^{1,0}$ and, by conjugation, to $H^{0,1}$.
Their intersection is zero.

On a finite root disk the local system is constant and $\psi$
extends as a nonzero real class to the central torus. Its
coherent image is nonzero there and invariant under the forward
deck. It supplies an invariant eigenframe of weight zero, so
its descent has no vanishing factor.
At the cusp, the neighborhood retracts onto $W$ and the
specialization kills the angular circles while preserving
the polygon periods. In the full adapted marking of
Section~\ref{hol:sec:adapted-fan}, these periods are $e_3,e_4$;
$\psi$ takes the values $0,1$ on them. Thus it extends to a
class on the cusp neighborhood whose restriction to $W$ is
nonzero. Inclusion $\mathbb C\to\mathcal O$ and the Leray
edge map give a holomorphic section across the disk.
By base change its central value is the nonzero image in
Lemma~\ref{nf:lem:cusp-coherent}. Functoriality identifies it
on the punctured disk with $\psi^{0,1}$. These extensions
give the subbundle $\mathcal O_B\eta\subset\mathcal G$.

We next determine $\mathcal C_2$. Equation~\eqref{nf:eq:canonical}
gives
\[
 \omega_{Y/B}=K_Y\otimes f^*K_B^{-1}
       \simeq\mathcal O_Y(2F_4)\otimes f^*\mathcal O_B(1).
\]
One has $f_*\mathcal O_Y(2F_4)=\mathcal O_B$.
For a local section, regard it as a meromorphic function
with poles bounded by $2F_4$. It is constant on each nearby
compact regular fiber. With $z=t-1$, multiplication by $z$
makes it holomorphic upstairs, so it descends to a holomorphic
function $b(z)$ because $f_*\mathcal O_Y=\mathcal O_B$.
The original function is $b(z)/z$ off the central fiber.
If $b(0)\ne0$, its pullback would have a pole of order four,
contrary to the bound two. Thus $b/z$ is holomorphic.
The projection formula now gives
$f_*\omega_{Y/B}=\mathcal O_B(1)$.

Proper analytic duality \cite{RamisRugetVerdier} reads here
\begin{equation}\label{nf:eq:proper-duality}
 Rf_*\omega_{Y/B}[2]
 \simeq R\mathcal Hom_B(Rf_*\mathcal O_Y,\mathcal O_B).
\end{equation}
The total space and base are smooth, and the cohomology
sheaves on the right are locally free by the preceding
argument. Its degree $-2$ sheaf is $\mathcal C_2^\vee$.
Hence $\mathcal C_2=\mathcal O_B(-1)$. This argument applies
also at the multiple and nonnormal fibers.

Cup product is an isomorphism
$\bigwedge^2\mathcal G\to\mathcal C_2$ on the smooth locus
and at the cusp by Lemma~\ref{nf:lem:cusp-coherent}.
At a finite point let $f_0=\eta$ and $f_1$ be root eigenframes.
The nontrivial holomorphic-form character from
\eqref{nf:eq:horizontal-frames} is $a$, with
\[
 (m,a)=(3,2)\quad\hbox{at }0,\qquad
 (m,a)=(4,1)\quad\hbox{at }1.
\]
On coherent $H^1$, the characters are $0,-a$; this follows
by conjugation on the central torus and then by the
equivariant eigenframe construction. The descended frames are
\begin{equation}\label{nf:eq:coherent-frames}
 \mathcal G:\ (f_0,s^af_1),\qquad
 \mathcal C_2:\ s^a(f_0\wedge f_1).
\end{equation}
Their cup product is exactly the displayed generator.
There is no additional base factor, since $a<m$.
Thus $\det\mathcal G=\mathcal C_2=\mathcal O_B(-1)$.
The section $\eta$ gives \eqref{nf:eq:G-sequence}.
It splits because
$\operatorname{Ext}^1(\mathcal O_B(-1),\mathcal O_B)
=H^1(B,\mathcal O_B(1))=0$.
Leray on the base curve gives the stated structure-sheaf
dimensions, and Serre duality gives those of $K_Y$.
\end{proof}

The section $\eta$ is a coherent higher-direct-image section.
It is not a global constant-coefficient lift of $\psi$.
The integral constructible extension in (5.4) is
$12\mathbb Z\psi$, and its nonzero transgression remains
part of the integral calculation.

\subsection{Higher direct images of relative forms}

Put
\[
 \mathcal A_q=R^qf_*\mathscr E,\qquad
 \mathcal B_q=R^qf_*\mathcal K_{\mathrm{rel}}.
\]
The scalar sequence supplies their global filtration.
The finite characters determine how its multiplication maps
extend across the two multiple fibers.

\begin{theorem}\label{nf:thm:relative-direct-images}
The higher direct images are locally free and vanish for $q>2$.
For $q=0,1,2$ their isomorphism types are
\begin{equation}\label{nf:eq:relative-direct-images}
 \begin{aligned}
 \mathcal B_0&=\mathcal O_B(-1),&
 \mathcal A_0&\simeq\mathcal O_B\oplus\mathcal O_B(-1),\\
 \mathcal B_1&\simeq\mathcal O_B(-1)\oplus\mathcal O_B,&
 \mathcal A_1&\simeq\mathcal O_B^{\oplus2}
                         \oplus\mathcal O_B(-1)^{\oplus2},\\
 \mathcal B_2&\simeq\mathcal O_B,&
 \mathcal A_2&\simeq\mathcal O_B\oplus\mathcal O_B(-1).
 \end{aligned}
\end{equation}
They carry the exact sequences induced by contraction,
\begin{equation}\label{nf:eq:A-filtration}
 0\longrightarrow\mathcal B_q
 \longrightarrow\mathcal A_q
 \longrightarrow\mathcal C_q\longrightarrow0.
\end{equation}
Multiplication by $\eta$ embeds
$\mathcal L=\mathcal B_0$ as a subbundle of $\mathcal B_1$,
and the quotient and degree-two line are
\begin{equation}\label{nf:eq:elementary-modifications}
 \mathcal B_1/\mathcal L
   \simeq\mathcal L\otimes\mathcal Q(0+1)\simeq\mathcal O_B,
 \qquad
 \mathcal B_2\simeq\mathcal L\otimes\mathcal C_2(0+1)
                   \simeq\mathcal O_B.
\end{equation}
In particular $h^q(Y,\mathscr E)=(1,2,1,0)$.
\end{theorem}

\begin{proof}
On a smooth torus family, evaluation identifies its relative
cotangent bundle with the pullback of its bundle of invariant
forms. This applies to the good-reduction root families.
On a root disk the scalar sequence has an equivariant splitting
by averaging. At the cusp the relative logarithmic toric frame
splits it directly. Exact finite invariants and the projection
formula consequently give local freeness and make all the
connecting maps of \eqref{nf:eq:scalar-sequence} zero.
They give the global sequences \eqref{nf:eq:A-filtration}.
Equivalently, the global form $\theta$ already supplies their
splitting with respect to contraction.

To compute the kernel terms, choose root frames $e_0,\alpha$
for relative forms and $f_0,f_1$ for coherent $H^1$, with
weights $(0,a)$ and $(0,-a)$, where $f_0=\eta$.
The invariant generator of the horizontal line is
$s^{m-a}\alpha$. Thus multiplication by $\eta$ gives the
primitive invariant generator
$s^{m-a}\alpha f_0$ of a subline of $\mathcal B_1$.
It is a subbundle on every chart.

Away from $0,1$, including at infinity, projection formula
identifies the quotient with $\mathcal L\otimes\mathcal Q$.
At a finite point the product of its descended generators is
\begin{equation}\label{nf:eq:one-base-factor}
 (s^{m-a}\alpha)(s^af_1)=s^m\alpha f_1
             =t_{\mathrm{loc}}\,\alpha f_1,
\end{equation}
where $t_{\mathrm{loc}}=t$ at $0$ and $t-1$ at $1$.
The invariant quotient is generated by $\alpha f_1$ itself.
Hence the product map vanishes to order exactly one at each
finite point and nowhere else. This proves the first formula
in \eqref{nf:eq:elementary-modifications}. The same calculation
with $f_0\wedge f_1$ proves the second. These identities specify
the lines through the divisors of their multiplication maps.

For clarity, the invariant frames of $\mathcal A_1$ are
\[
 e_0f_0,\quad s^ae_0f_1,\quad
 s^{m-a}\alpha f_0,\quad\alpha f_1.
\]
In these frames the natural maps are
\begin{equation}\label{nf:eq:multiplication-matrices}
 \mathcal H\otimes\mathcal G\longrightarrow\mathcal A_1:
       \operatorname{diag}(1,1,1,t_{\mathrm{loc}}),
 \qquad
 \mathcal H\otimes\mathcal C_2\longrightarrow\mathcal A_2:
       \operatorname{diag}(1,t_{\mathrm{loc}}).
\end{equation}
This exhibits the two elementary modifications without
replacing the actual invariant modules by their generic fibers.

Now $\mathcal B_0=\mathcal L=\mathcal O_B(-1)$, and
\[
 0\longrightarrow\mathcal O_B(-1)
 \longrightarrow\mathcal B_1\longrightarrow\mathcal O_B
 \longrightarrow0
\]
splits because $H^1(\mathcal O_B(-1))=0$.
Insert these kernel terms and \eqref{nf:eq:coherent-bundles}
into \eqref{nf:eq:A-filtration}. Every extension between
$\mathcal O_B$ and $\mathcal O_B(-1)$ splits, since the
relevant $H^1(\mathcal O_B(k))$ has $k\ge-1$.
This proves \eqref{nf:eq:relative-direct-images}.
The specified filtration and multiplication maps are retained
even when the pure line-summand decompositions are noncanonical.
Leray gives the last assertion.
\end{proof}

\subsection{The Picard group and tangent-cohomology lengths}

The coherent section $\eta$ also determines which line bundles
on $Y$ restrict nontrivially to a chosen smooth torus fiber.

\begin{corollary}\label{nf:cor:picard}
There is an additive isomorphism
$\operatorname{Pic}(Y)\simeq\mathbb C$.
Choose its coordinate $\lambda$ so that the corresponding
class in $H^1(Y,\mathcal O_Y)$ maps to $\lambda\eta$, and
write $L_\lambda$ for its exponential line bundle.
For every fixed smooth fiber $F$, all but countably many
$\lambda$ have non-torsion restriction $L_\lambda|_F$.
For any such $\lambda$, the sheaves
$R^if_*(T_Y\otimes L_\lambda)$, $i=1,2$, have finite support,
and their lengths satisfy
\[
 \ell_2=\ell_1+1,\qquad \ell_2\ge1.
\]
\end{corollary}

\begin{proof}
Theorems~\ref{cmp:thm:comparison} and
\ref{th:thm:real-sphere} give
$H^1(Y;\mathbb Z)=H^2(Y;\mathbb Z)=0$.
The exponential sequence, with exponential
$h\mapsto\exp(2\pi i h)$, therefore identifies
$\operatorname{Pic}(Y)$ with $H^1(Y,\mathcal O_Y)$.
The coherent Leray edge map is an isomorphism
\[
 H^1(Y,\mathcal O_Y)\xrightarrow{\ \sim\ }
 H^0(B,\mathcal G)=\mathbb C\eta,
\]
since $H^1(B,\mathcal O_B)=H^2(B,\mathcal O_B)=0$.

For a fixed smooth $F$, naturality of the edge and base-change
maps identifies the restriction of $L_\lambda$ with
\[
 \lambda\eta_F\pmod{\Lambda_F},\qquad
 \Lambda_F=\operatorname{im}\bigl(
 H^1(F;\mathbb Z)\longrightarrow H^1(F,\mathcal O_F)\bigr).
\]
The projection $H^1(F;\mathbb R)\to H^{0,1}(F)$ is a
real-linear isomorphism, by the Hodge decomposition of the
complex torus and conjugation. Thus $\Lambda_F$ is a full
lattice, and $\eta_F\ne0$. Torsion means that
$n\lambda\eta_F\in\Lambda_F$ for some integer $n>0$.
For each $n$ and lattice vector there is at most one
$\lambda$, so the exceptional set is countable.

Choose $\lambda$ outside it. Its restriction is nontrivial
and topologically trivial, and
Proposition~\ref{ad:prop:one-fiber-lengths} applies.
Proposition~\ref{ad:prop:reduction} gives
$\chi(Y,T_Y\otimes L_\lambda)=1$, hence
$\ell_2-\ell_1=1$. Corollary~\ref{ad:cor:ambient-nonvanishing}
also gives a nonzero fiber of the second direct image at
infinity, so $\ell_2\ge1$.
\end{proof}

This argument fixes one smooth fiber. It gives neither a
simultaneous non-torsion restriction on all fibers nor the
individual lengths or local elementary divisors. Its coherent
edge isomorphism is compatible with the nonzero integral
transgression of $12\psi$; it does not replace that
transgression by a constant-coefficient lift.

\section{Ordinary Hodge groups and the completed de Rham complex}
\label{sec:hodge-rees}

Relative logarithmic forms retain two further maps. Their
extension by base differentials gives the Kodaira--Spencer
maps, and their branch residues compare them with ordinary
forms. We compute these maps before passing to cohomology.

\subsection{The logarithmic Kodaira--Spencer maps}

Applying $Rf_*$ to \eqref{nf:eq:absolute-log-sequence} gives
\begin{equation}\label{hr:eq:KS-maps}
 \kappa_q:\mathcal A_q\longrightarrow
                 \mathscr L_B\otimes\mathcal C_{q+1}.
\end{equation}
The base arrow is the actual pullback of logarithmic
differentials, including the fiber multiplicities.
Recall the specified subline
$\mathcal O_B\eta\subset\mathcal G$, its quotient $\mathcal Q$,
and the unique relative form $\theta$ with $\theta(V)=1$.

\begin{proposition}\label{hr:prop:KS}
Let $\mathcal T=\mathbb C_0\oplus\mathbb C_1$ be the sum of
the two residue-field sheaves on $B$. There is an exact sequence
\begin{equation}\label{hr:eq:KS-zero}
 0\longrightarrow\mathcal H\xrightarrow{\ \kappa_0\ }
 \mathscr L_B\otimes\mathcal G
 \longrightarrow\mathcal T\longrightarrow0.
\end{equation}
The map
$\kappa_1:\mathcal A_1\to\mathscr L_B\otimes\mathcal C_2
\simeq\mathcal O_B$ is onto both as a sheaf map and on global
sections. Its kernel is isomorphic to
$\mathcal O_B\oplus\mathcal O_B(-1)^{\oplus2}$.
\end{proposition}

\begin{proof}
At the cusp use its actual adapted period presentation
\[
 (z,w)\sim(qc_1(q)z,c_2(q)w),\qquad
 (z,w)\sim(c_3(q)z,qc_4(q)w),
\]
with all four $c_i$ holomorphic units.
The absolute lifts $d\log z,d\log w$ of the relative frame
have deck differences, under the period $(k_1,k_2)$,
\[
 \left((k_1,k_2)+q\,h_k(q)\right)d\log q
\]
for a holomorphic vector-valued function $h_k$.
At $q=0$ these are precisely the two polygon-period
covectors. Lemma~\ref{nf:lem:cusp-coherent} identifies their
coherent images with a basis. Thus $\kappa_0$ is invertible
at infinity. The simultaneous coordinate and period shear
changes the two bases by invertible integral matrices.
The unit derivatives are retained in $q\,h_k(q)$ and
vanish at the central value.

At a finite point use $t_{\mathrm{loc}}=s^m$ and the forward
weights $(0,a)$ for root relative forms $e_0,e_1$, with
$e_1$ horizontal, and $(0,-a)$ for coherent frames $f_0,f_1$,
with $f_0=\eta$. Here $(m,a)=(3,2)$ or $(4,1)$.
The descended frames are
\[
 \mathcal H:\ (e_0,s^{m-a}e_1),\qquad
 \mathcal G:\ (f_0,s^af_1).
\]
Write $k_{ji}(s)$ for the raw matrix of the pulled-back
$\kappa_0$, using the base generator
$dt_{\mathrm{loc}}/t_{\mathrm{loc}}=m\,ds/s$.
The ordinary smooth-family connecting map has a holomorphic
$ds$ coefficient, so every $k_{ji}$ is divisible by $s$.
The factor $1/m$ is included in this convention.
Equivariance says
\begin{equation}\label{hr:eq:KS-weights}
 k_{ji}(\zeta s)=\zeta^{\,b_i-c_j}k_{ji}(s),
 \qquad b=(0,a),\quad c=(0,-a).
\end{equation}
For residue class zero, the least allowed exponent is $m$,
since the constant term is forbidden.

Multiplication by the source exponents and division by the
target exponents now give
\begin{equation}\label{hr:eq:finite-KS-matrices}
 M_3(t_{\mathrm{loc}})
  =\begin{pmatrix}t_{\mathrm{loc}}A&t_{\mathrm{loc}}B\\ C&D\end{pmatrix},
 \qquad
 M_4(t_{\mathrm{loc}})
  =\begin{pmatrix}t_{\mathrm{loc}}A&t_{\mathrm{loc}}B\\
                       C&t_{\mathrm{loc}}D\end{pmatrix}.
\end{equation}
$A,B,C,D$ denote holomorphic functions, separately on
the two disks. For example, the order-four raw entry
$k_{11}$ starts at $s^2$. Its source factor is $s^3$
and its target factor $s$, giving $s^4=t_{\mathrm{loc}}$.
This explains the lower-right vanishing in $M_4$.

The global determinant of $\kappa_0$ is a section of
\[
 \det(\mathscr L_B\otimes\mathcal G)\otimes
       (\det\mathcal H)^{-1}\simeq\mathcal O_B(2).
\]
It is nonzero at infinity and vanishes at least once at
each finite point by \eqref{hr:eq:finite-KS-matrices}.
Those two zeros are therefore simple and there are no
others. Each finite matrix has rank one at its center.
Its first row is zero, so its image is the $f_1$ line.
The surviving quotient is the $f_0$ direction, proving
\eqref{hr:eq:KS-zero}.

On the smooth locus the next connecting map is cup
product with the first:
\begin{equation}\label{hr:eq:KS-cup}
 \kappa_1(hv)=\kappa_0(h)\wedge v.
\end{equation}
The same identity holds globally on the image of the
actual multiplication
$\mathcal H\otimes\mathcal G\to\mathcal A_1$.
Multiplication by $f_0=\eta$ is unmodified at the finite
points by \eqref{nf:eq:multiplication-matrices}.
Their nonzero $f_1$ images therefore wedge nontrivially
with $f_0$. Elsewhere $\kappa_0$ is invertible.
This proves sheaf surjectivity of $\kappa_1$.

To prove surjectivity on global sections, compose
$\kappa_0$ with the specified quotient:
\[
 b:\mathcal H\longrightarrow
       \mathscr L_B\otimes\mathcal Q\simeq\mathcal O_B.
\]
At the order-four point the lower-right entry in $M_4$
shows that $b$ vanishes on the horizontal fiber.
On the other hand,
\[
 \det M_4=t_{\mathrm{loc}}
            (t_{\mathrm{loc}}AD-BC)
\]
has a simple zero. Thus $B(0)C(0)\ne0$, and $b$ is
nonzero on the complementary direction.
Since $\theta$ has quotient value one in
$\mathcal H/\mathcal L$, $b(\theta)$ is nonzero there.
It is a global section of a degree-zero line and hence
a nonzero constant. Equation~\eqref{hr:eq:KS-cup} makes
$\kappa_1(\theta\eta)$ nonzero. After a scalar
normalization this is a right inverse
$\mathcal O_B\to\mathcal A_1$.
Subtracting its multiples from the other summands of
\eqref{nf:eq:relative-direct-images} identifies the kernel
with $\mathcal O_B\oplus\mathcal O_B(-1)^{\oplus2}$.
\end{proof}

The nonvanishing just proved uses the order-four matrix
and the global quotient line. It requires no identification
with an auxiliary normal function or its derivative.

\subsection{Absolute logarithmic forms and normalized residues}

Write
$\mathcal P_q=R^qf_*\Omega_Y^1(\log D)$ and
$\mathcal K_1=\ker\kappa_1$.
The long exact sequence of
\eqref{nf:eq:absolute-log-sequence}, together with
Proposition~\ref{hr:prop:KS}, gives
\begin{equation}\label{hr:eq:absolute-direct-images}
 \mathcal P_0=\mathscr L_B,\qquad
 0\longrightarrow\mathcal T\longrightarrow\mathcal P_1
 \longrightarrow\mathcal K_1\longrightarrow0,\qquad
 \mathcal P_2=\mathcal A_2.
\end{equation}
The middle extension splits because
$H^1(B,\mathcal K_1^\vee\otimes\mathcal T)=0$.
Its specified torsion inclusion remains the one induced
by $\mathscr L_B\otimes\mathcal G$.
Consequently
\[
 \mathcal P_1\simeq
 \mathcal O_B\oplus\mathcal O_B(-1)^{\oplus2}
                \oplus\mathbb C_0\oplus\mathbb C_1,
 \qquad
 h^q\bigl(Y,\Omega_Y^1(\log D)\bigr)=(2,3,1,0).
\]
In particular, every global absolute logarithmic
one-form is pulled back from $B$.

The normalization of $D$ is
$\widetilde D=S\sqcup F_3\sqcup F_4$.
Let $\nu_D:\widetilde D\to D\subset Y$ be its map and
put $\mathcal R_D=(\nu_D)_*\mathcal O_{\widetilde D}$.
The residue sequence is
\begin{equation}\label{hr:eq:residue-sequence}
 0\longrightarrow\Omega_Y^1
 \longrightarrow\Omega_Y^1(\log D)
 \xrightarrow{\ \operatorname{Res}\ }\mathcal R_D
 \longrightarrow0.
\end{equation}
Indeed, on a local branch chart the residue of
$\sum h_i\,dx_i/x_i$ is $h_i|_{x_i=0}$ on each branch.
Each branch function has a holomorphic lift locally;
the kernel consists of ordinary holomorphic forms.
These maps descend through branch permutations.

\begin{theorem}\label{hr:thm:hodge}
The ordinary one-form cohomology is
\[
 h^q(Y,\Omega_Y^1)=(0,2,1,0).
\]
With $h^{p,q}(Y)=\dim H^q(Y,\Omega_Y^p)$, the full table is
\begin{equation}\label{hr:eq:hodge-table}
 \begin{array}{c|rrrr}
 q\backslash p&0&1&2&3\\ \hline
 0&1&0&0&0\\
 1&1&2&1&0\\
 2&0&1&2&1\\
 3&0&0&0&1
 \end{array}.
\end{equation}
\end{theorem}

\begin{proof}
The normalized toric surface $S$ has coherent dimensions
$(1,0,0)$. Each $F_m$ is the free finite quotient of its
central smooth torus. Taking invariants of the coherent
exterior algebra with weights $(0,-a)$ gives dimensions
$(1,1,0)$, since $a\not\equiv0\pmod m$.
Therefore
\[
 h^q(Y,\mathcal R_D)=(3,2,0).
\]

On global sections, residues are the pullback
base-residue map
\begin{equation}\label{hr:eq:base-residues}
 \beta\longmapsto
 \bigl(\operatorname{res}_\infty\beta,\,
       3\operatorname{res}_0\beta,\,
       4\operatorname{res}_1\beta\bigr).
\end{equation}
Base residues sum to zero and give a rank-two map.
The factors $3,4$ are the scheme-fiber multiplicities;
derivatives of units add no residue. Thus this map is
injective on the two-dimensional space of global
logarithmic one-forms.

We also need the map in degree one. At a finite point
let $\beta_m$ generate $H^1(F_m,\mathcal O_{F_m})$ so
that its pullback is the invariant central class $f_0$.
Extend $f_0$ by $\eta$ on the root disk. The class
\begin{equation}\label{hr:eq:normalized-residue}
 \frac1m\,\frac{dt_{\mathrm{loc}}}{t_{\mathrm{loc}}}
                       \otimes\eta
\end{equation}
maps into $\mathcal P_1$ and has residue $\beta_m$,
because its pullback is $(ds/s)\otimes\eta$.
The other descended coherent generator $s^af_1$
restricts to zero on the reduced central fiber.
The fiber image of $\kappa_0$ is the $f_1$ direction,
so residue identifies its quotient
$\mathcal T$ with $R^1f_*\mathcal R_D$ at each finite
point. At infinity the target is
$H^1(S,\mathcal O_S)=0$. Hence
\[
 \mathcal T\xrightarrow{\ \sim\ }R^1f_*\mathcal R_D.
\]

The Leray maps identify
$H^1(Y,\Omega_Y^1(\log D))$ with
$H^0(B,\mathcal P_1)$, since
$H^1(B,\mathcal P_0)=H^1(B,\mathcal O_B(1))=0$.
They identify $H^1(Y,\mathcal R_D)$ with
$H^0(B,R^1f_*\mathcal R_D)$, since $f_*\mathcal R_D$
has finite support. The torsion inclusion in
\eqref{hr:eq:absolute-direct-images} therefore proves
surjectivity of the global degree-one residue map.

In the long exact sequence of
\eqref{hr:eq:residue-sequence}, the degree-zero
cokernel has dimension one and the degree-one kernel
has dimension $3-2=1$. Thus $h^1(\Omega_Y^1)=2$.
Surjectivity in degree one and $H^2(\mathcal R_D)=0$
identify $H^2(\Omega_Y^1)$ with the one-dimensional
$H^2(\Omega_Y^1(\log D))$. The end terms give
$H^0(\Omega_Y^1)=H^3(\Omega_Y^1)=0$.
Combine this with Theorem~\ref{nf:thm:coherent} and
Serre duality to obtain \eqref{hr:eq:hodge-table}.
\end{proof}

\subsection{The Fr\"olicher differential}

We now also use the smooth comparison for this selected
threefold. Theorems~\ref{cmp:thm:comparison} and
\ref{th:thm:real-sphere} give
\begin{equation}\label{hr:eq:betti}
 b_0(Y)=b_6(Y)=1,\qquad b_k(Y)=0\quad(1\le k\le5).
\end{equation}
This topological input enters the following spectral-sequence
argument; the coherent sheaves above were computed on the
complex quotient itself.

\begin{proposition}\label{hr:prop:frolicher}
The Fr\"olicher spectral sequence of $Y$ degenerates at
$E_2$, whose only nonzero terms are $\mathbb C$ in
bidegrees $(0,0)$ and $(3,3)$. Its four nonzero $d_1$
arrows each have rank one.
\end{proposition}

\begin{proof}
Apart from the two endpoints, the $E_1$ page
\eqref{hr:eq:hodge-table} has two rows:
\[
 \begin{array}{ccl}
 q=1:&
 \mathbb C_{(0,1)}\longrightarrow
 \mathbb C^2_{(1,1)}\longrightarrow
 \mathbb C_{(2,1)},\\[2pt]
 q=2:&
 \mathbb C_{(1,2)}\longrightarrow
 \mathbb C^2_{(2,2)}\longrightarrow
 \mathbb C_{(3,2)}.
 \end{array}
\]
Here the subscripts record bidegrees and $d_1$ is induced
by the actual holomorphic differential $\partial$.
The term $E_2^{0,1}$ cannot be changed by a higher
differential: its only possible outgoing $d_2$ targets
$E_2^{2,0}=0$, and there is no incoming source.
Thus $b_1=0$ forces the first arrow to be injective.
The term $E_2^{1,1}$ has no higher source or target, so
$b_2=0$ forces the next arrow to have rank one.
Likewise $E_2^{1,2}$ and $E_2^{2,2}$ have no nonzero
higher source or target; $b_3=b_4=0$ gives rank one
for the two arrows in the second row.
Since $d_1^2=0$, both rows are exact.
Only the endpoints remain on $E_2$.
\end{proof}

\subsection{A finite model over the formal parameter}

Let $R=\mathbb C[[\hbar]]$, where $\hbar$ has cochain
degree zero. This formal parameter is distinct from
the contraction parameter $u$. Let
$\mathcal A_Y^{p,q}$ be the sheaf of smooth complex
$(p,q)$-forms and define the total cochain complex
\begin{equation}\label{hr:eq:Rees-definition}
 \mathcal R_Y=
 \left(\mathcal A_Y^{\bullet,\bullet}[[\hbar]],
                  d_\hbar=\bar\partial+\hbar\partial\right),
 \qquad
 K_{\mathrm{Rees}}(Y)=R\Gamma(Y,\mathcal R_Y).
\end{equation}
The completion is $\hbar$-adic. Its sheaves are fine,
with partitions of unity acting coefficientwise, so global
smooth forms represent the displayed derived sections.

\begin{theorem}\label{hr:thm:Rees-model}
The completed complex $K_{\mathrm{Rees}}(Y)$ admits a
finite free model with two zero-differential summands
$R$ in degrees $0,6$ and four summands
\begin{equation}\label{hr:eq:four-pairs}
 [\,Ra_j\xrightarrow{\ \hbar\ }Rb_{j+1}\,],
 \qquad |a_j|=j,\quad |b_{j+1}|=j+1,\quad 1\le j\le4.
\end{equation}
It is connected to the actual completed form complex
by the inclusion, projection and homotopy constructed
below.
\end{theorem}

\begin{proof}
Choose a Hermitian metric on compact $Y$.
Let $H$ be the finite-dimensional space of
$\bar\partial$-harmonic forms, $i:H\to A$ its inclusion
into $A=\Gamma(Y,\mathcal A_Y^{\bullet,\bullet})$,
and $p:A\to H$ harmonic projection.
Let $G_{\bar\partial}$ be the Green operator of the
Dolbeault Laplacian and put
$K=\bar\partial^*G_{\bar\partial}$.
The compact Hermitian elliptic theory
\cite[Part~I, Theorems~3.15 and 7.1--7.3]{DemaillyHodge}
gives
\begin{equation}\label{hr:eq:harmonic-contraction}
 \bar\partial K+K\bar\partial=1-ip,\qquad
 K^2=0,\quad Ki=0,\quad pK=0,\quad pi=1.
\end{equation}
Only the Dolbeault Laplacian is used; no K\"ahler
identity is required.

Put $\delta=\hbar\partial$ and define
\begin{equation}\label{hr:eq:perturbation}
 \begin{aligned}
 i'&=(1+K\delta)^{-1}i,&
 p'&=p(1+\delta K)^{-1},\\
 K'&=K(1+\delta K)^{-1},&
 D_H&=p\delta(1+K\delta)^{-1}i.
 \end{aligned}
\end{equation}
Each inverse is a finite geometric series on forms:
$K\delta$ and $\delta K$ increase holomorphic degree
by one, and that degree is at most three.
Using \eqref{hr:eq:harmonic-contraction} and
$\bar\partial\delta+\delta\bar\partial=\delta^2=0$,
substitution gives
\begin{equation}\label{hr:eq:perturbed-identities}
 \begin{gathered}
 d_\hbar i'=i'D_H,\qquad p'd_\hbar=D_Hp',\qquad p'i'=1,\\
 d_\hbar K'+K'd_\hbar=1-i'p',\qquad D_H^2=0.
 \end{gathered}
\end{equation}
For example, $(1+K\delta)i'=i$ and
$(1+\delta K)^{-1}\delta
=\delta(1+K\delta)^{-1}$ reduce the first identity
to \eqref{hr:eq:harmonic-contraction}; the remaining
identities follow by the same finite cancellation.
Thus these are actual comparison maps and a homotopy.

The transferred differential has the form
\[
 D_H=\hbar d_1+\hbar^2d_{(2)}+\hbar^3d_{(3)},\qquad
 d_1=p\partial i,
\]
where
$d_{(r)}=(-1)^{r-1}p\partial(K\partial)^{r-1}i$
has bidegree $(r,1-r)$.
For $r\ge2$, table~\eqref{hr:eq:hodge-table} has no
pair of nonzero source and target in that bidegree.
Hence
\begin{equation}\label{hr:eq:transferred-differential}
 D_H=\hbar d_1
\end{equation}
exactly. This conclusion uses the bidegrees of the
transferred operators, in addition to
Proposition~\ref{hr:prop:frolicher}.

Choose $a_1\in H^{0,1}$ and $a_3\in H^{1,2}$ with
nonzero $d_1$ images. Choose $a_2\in H^{1,1}$ in a
complement to $d_1H^{0,1}$, and $a_4\in H^{2,2}$
in a complement to $d_1H^{1,2}$, again with nonzero
images. Set $b_{j+1}=d_1a_j$.
The two exact rows show that these eight vectors,
together with $1$ and a normalized top harmonic
class, form a basis of $H$.
Their images are defined by the actual operator
$p\partial i$. Extending this basis to $R$ gives
\eqref{hr:eq:four-pairs}, with the comparison maps
\eqref{hr:eq:perturbation} unchanged.
\end{proof}

\subsection{Cohomology and the derived special fiber}

The model makes the formal torsion and its specialization
explicit:
\begin{equation}\label{hr:eq:Rees-cohomology}
 H^n(K_{\mathrm{Rees}}(Y))=
 \begin{cases}
 R,&n=0,6,\\
 R/(\hbar),&n=2,3,4,5,\\
 0,&n=1\text{ or }n\notin[0,6].
 \end{cases}
\end{equation}
Every torsion summand has order exactly one.
After inverting $\hbar$, multiplication of a
$(p,q)$-form by $\hbar^{-p}$ intertwines
$\bar\partial+\hbar\partial$ with
$\bar\partial+\partial$. Since $p$ is bounded, this
is a well-defined isomorphism of Laurent-series
complexes. It identifies the generic fiber with
ordinary de Rham cohomology over $\mathbb C((\hbar))$,
with its two free classes in degrees zero and six.

Reduction of the finite free model modulo $\hbar$
computes its derived special fiber. Its differential
is zero, and the total-degree dimensions are
\begin{equation}\label{hr:eq:derived-dimensions}
 \dim_{\mathbb C}
 H^n\!\left(K_{\mathrm{Rees}}(Y)\otimes_R^{\mathbb L}\mathbb C\right)
       =(1,1,2,2,2,1,1)\quad(0\le n\le6).
\end{equation}
These are the total Dolbeault dimensions.
The ordinary tensor product of
\eqref{hr:eq:Rees-cohomology} instead has dimensions
$(1,0,1,1,1,1,1)$.
The difference is expressed by
\begin{equation}\label{hr:eq:derived-Tor}
 \begin{split}
 0\longrightarrow&
 H^n(K_{\mathrm{Rees}}(Y))\otimes_R\mathbb C\\
 \longrightarrow&
 H^n\!\left(K_{\mathrm{Rees}}(Y)
                    \otimes_R^{\mathbb L}\mathbb C\right)\\
 \longrightarrow&
 \operatorname{Tor}_1^R
       \bigl(H^{n+1}(K_{\mathrm{Rees}}(Y)),\mathbb C\bigr)
 \longrightarrow0.
 \end{split}
\end{equation}
For the pair indexed by $j$, $b_{j+1}$ represents
torsion cohomology in degree $j+1$. The vector $a_j$
is not a formal cycle lifting its Dolbeault class.
It survives in the derived special fiber as the
corresponding $\operatorname{Tor}_1$ class in degree
$j$. Thus the four additional contributions occur
in degrees $1,2,3,4$.

\subsection{Products, unit and trace}

The completed form complex has its wedge product and
constant unit. The Laurent-series comparison
$\alpha^{p,q}\mapsto\hbar^{-p}\alpha^{p,q}$ preserves
both. The finite harmonic contraction retains these
operations by homotopy transfer. Its binary operation is
\[
 m_2(a,b)=p'\bigl(i'a\wedge i'b\bigr).
\]
The higher operations are the finite sums over planar
binary trees with $i'$ at the leaves, wedge product
at each vertex, $K'$ on internal edges and $p'$ at
the root. The signs are the cochain Koszul signs,
with $K'$ of degree $-1$; equivalently the same
formula on the suspended tensor coalgebra uses the
degree-one coderivations. The side conditions in
\eqref{hr:eq:harmonic-contraction} give the transferred
unital $A_\infty$ structure and its higher comparison
maps. The direct-sum basis in
\eqref{hr:eq:four-pairs} describes the differential;
it does not assert strict preservation of wedge
multiplication in that basis.

Integration of the total-degree-six form $i'a$
defines the trace. Stokes' theorem and
\eqref{hr:eq:perturbed-identities} make it vanish on
$D_H$-boundaries. The constant unit and the top
harmonic class have no perturbation correction:
$\partial$ annihilates both. Normalize the top class
to have integral one.

All constructions above use the fixed admissible
$u$ and its specified lifts and adapted gluing.
They apply to each choice satisfying these hypotheses,
but do not construct descent around the parameter
monodromy or a filling at $u=0$.
The completed complex \eqref{hr:eq:Rees-definition}
is a cochain complex of forms, with a formal parameter;
it supplies no loop quantization or chiral realization.
The actual double complex determines the Bott--Chern and
Aeppli groups in Section~\ref{sec:bott-chern}.
Section~\ref{sec:courant} computes the Aeppli class map for
pluriclosed forms. The same logarithmic residue sequence will
give the coherent tangent comparison in
Section~\ref{sec:tangent-curves}, including its double curves
and triple points. The actual residue maps also lift the
invariant integral fiber class to ordinary Dolbeault cohomology.
Its cusp restriction evaluates the curve periods and excludes
pluriclosed Hermitian metrics on this selected threefold.
Comparison with a different marked rank-four family still
requires its own geometric maps.
